\section{Implementation}\label{sec:impl}

% One subsection per major code surface; ~1 paragraph each.
% Code references should point at github SHAs of the open-source
% release. PHASES.md tags each phase with its core commit.

\subsection{Coordinator and worker runtime}
% - Python 3.10, gRPC, PyTorch (Jetson-specific NV CUDA wheels)
% - 99 unit/smoke tests, 4 of which exercise live multi-node behavior
% - Ansible deploy for the 5-host Jetson fleet
% Numbers: lines of code, # of RPCs in radp.proto, etc.

\subsection{gRPC protocol}
% - Table of RPCs: LoadStage, RunStage (+ async_chain field),
%   MirrorActivation, ResultReady, SetNextHop, LoadHead,
%   Heartbeat, EvictRequest, PromoteBackup, ProfileLayers,
%   MeasurePeer, Ping
% - Trailer metadata scheme (radp-failed-start, radp-failed-end)

\subsection{Auto-scheduling}
% - Coord profiles every worker (per-layer compute, peer-to-peer
%   network) on boot then runs the DP. The profile sidecar (
%   /tmp/radp_scheduler_stats.json) is the single source of truth
%   for the experiments in \S\ref{sec:eval}.

\subsection{Failure detection and recovery}
% - HeartbeatSender + FailureDetector
% - When trailer-based attribution beats heartbeat: the recovery
%   driver's "already-dead" branch finishes whatever heartbeat
%   started (rewire + replay) without redundant mark_dead
% - When the gateway's async timeout beats heartbeat: head
%   fallback (no trailer to read), accuracy preserved by R+psi
%   feasibility, efficiency degraded

\todo{Fill each subsection. Cross-reference \texttt{PHASES.md}
sections inline so reviewers can audit any specific implementation
claim against the live measurement that backed it.}
